\documentclass[10pt,twocolumn]{article}

% --- Standard conference/journal-like packages ---
\usepackage[margin=0.75in]{geometry}
\usepackage{times}
\usepackage{mathptmx}
\usepackage{amsmath,amssymb,amsfonts}
\usepackage{booktabs}
\usepackage{graphicx}
\usepackage{microtype}
\usepackage{natbib}
\usepackage{hyperref}
\usepackage{url}
\usepackage{enumitem}
\usepackage{caption}
\usepackage{subcaption}

% --- Tighten spacing to look more like an ML paper ---
\setlength{\columnsep}{0.25in}
\setlength{\textfloatsep}{8pt plus 2pt minus 2pt}
\setlength{\abovecaptionskip}{4pt}
\setlength{\belowcaptionskip}{0pt}
\setlist[itemize]{leftmargin=1.2em,topsep=2pt,itemsep=1pt,parsep=0pt}
\setlist[enumerate]{leftmargin=1.2em,topsep=2pt,itemsep=1pt,parsep=0pt}

\hypersetup{
  colorlinks=true,
  linkcolor=blue,
  citecolor=blue,
  urlcolor=blue
}

\title{Angular Manifold Routing: Fixed Geometric Routing from Embedding Angular Structure}

\author{
Casey Allard\\
Independent Researcher
}

\date{}

\begin{document}
\maketitle

\begin{abstract}
We study whether the angular non-uniformity of normalized token embeddings can be exploited for routing computation in transformer-style architectures. We propose a fixed geometric routing method based on a Hopf-fibration-derived sector map, and evaluate it in two settings: a semantic proxy based on PPMI-SVD embeddings, and a small trainable language model.

On the semantic proxy, the effective routing footprint grows sublinearly with the routing budget, following an empirical scaling law of approximately $K^{0.572}$, compared with $K^{1.0}$ for dense routing. This advantage persists through $K=5000$. In a 2-layer trainable language model, fixed geometric routing replaces a learned top-1 gate with an 8\% validation perplexity cost and no learned gate matrix, while retaining a concentrated routing footprint. A second-dataset replication on WikiText-2 matches the Penn Treebank result closely.

A key limitation is that in the trainable language model, the specific Hopf sector geometry is not distinguishable from randomly permuted fixed routing, suggesting that expert weights co-adapt to any fixed routing scaffold in this small-scale regime. The contribution of this work is therefore narrow: angular structure supports sparse fixed routing and transfers partially to trainable routing, but geometry-specific quality gains are not established in end-to-end language model training.
\end{abstract}

\section{Introduction}
Routing and sparsity are central themes in efficient transformer design. Mixture-of-experts models and related sparse architectures improve scaling by activating only a subset of available compute paths for a given input. Most such approaches rely on learned gating, which introduces additional parameters and training dynamics of its own.

This paper investigates a different question: can some routing structure be obtained directly from the geometry of the representation space, without a learned routing network? The motivating observation is that normalized token embeddings are not uniformly distributed on the unit hypersphere. If this angular structure is sufficiently stable, then a fixed geometric partition of the sphere may induce non-uniform routing that concentrates traffic into a subset of available sectors.

We study this question using a fixed routing map derived from a 4D Hopf fibration projection. The paper makes three scoped contributions:
\begin{enumerate}
    \item It shows that on a semantic proxy, fixed geometric routing produces a sublinear effective routing footprint over a broad range of routing budgets.
    \item It shows that this routing mechanism transfers partially to a small trainable language model, where it replaces a learned top-1 gate at modest perplexity cost.
    \item It shows that in this trainable setting, the specific Hopf geometry does not outperform randomly permuted fixed routing, clarifying the scope of the result.
\end{enumerate}

The intended claim is narrow. This paper does not establish broad replacement of learned MoE routing in realistic large-scale systems. It presents an empirical efficiency result in a proxy setting and a toy-scale trainable language model.

\section{Background}
Recent work on compression has shown that normalized language-model embeddings exhibit exploitable angular structure. This motivates asking whether the same structure can be used for routing rather than quantization.

Our focus here is not compression, but fixed routing. The central hypothesis is that angular non-uniformity may support a sparse routing map with no learned gate matrix.

\section{Method}
\subsection{Fixed geometric routing}
Given an L2-normalized input vector $x \in \mathbb{R}^D$, we project to a 4D subspace using four fixed coordinates:
\[
(a,b,c,d) = (x_i, x_j, x_k, x_l).
\]
We then compute Hopf-base coordinates:
\[
\rho_1 = \sqrt{a^2+b^2}, \qquad \rho_2 = \sqrt{c^2+d^2},
\]
\[
\chi_u = \frac{\rho_2^2}{\rho_1^2+\rho_2^2}, \quad
\delta = \operatorname{atan2}(b,a) - \operatorname{atan2}(d,c),
\]
\[
\alpha = \frac{1}{2}\left(\operatorname{atan2}(b,a)+\operatorname{atan2}(d,c)\right).
\]
An optional transport correction gives
\[
\alpha_t = \alpha + \frac{1}{2}\lambda \cos(2\chi)\,\delta.
\]
These coordinates are discretized into a product sector budget $K$, yielding a fixed routing sector.

\subsection{Effective routing footprint}
To measure routing concentration, we use an effective-buckets metric $\hat e$, which summarizes how many sectors are effectively active under the empirical routing distribution. Lower $\hat e$ relative to $K$ indicates stronger concentration and lower effective routing footprint.

\section{Proxy Experiments}
\subsection{Setup}
We first evaluate the routing mechanism on a semantic proxy based on PPMI-SVD embeddings. This setting isolates routing geometry from the confound of end-to-end expert co-adaptation.

We compare the fixed Hopf routing map against dense routing and randomized controls, and sweep routing budgets from small to large $K$.

\subsection{Main result}
Across the proxy experiments, the effective routing footprint scales approximately as
\[
\hat e \approx 2.957 \cdot K^{0.572}
\]
for the canonical Hopf routing configuration, compared with linear $K^{1.0}$ scaling for dense routing.

At larger budgets, the advantage persists through $K=5000$, with observed ratios in the range of approximately $2.6$ to $2.8\times$.

\begin{table}[t]
\centering
\caption{Proxy routing scaling laws.}
\label{tab:scaling}
\small
\begin{tabular}{lcc}
\toprule
Routing & Scaling law & $R^2$ \\
\midrule
HOPF & $\hat e = 2.957 \cdot K^{0.572}$ & 0.963 \\
PERM & $\hat e = 1.664 \cdot K^{0.814}$ & 0.993 \\
DENSE & $\hat e = K^{1.0}$ & 1.000 \\
\bottomrule
\end{tabular}
\end{table}

\begin{figure}[t]
    \centering
    \includegraphics[width=\columnwidth]{figures/fig_scaling_law.pdf}
    \caption{Effective routing footprint versus routing budget. HOPF shows sublinear growth relative to dense routing.}
    \label{fig:scaling}
\end{figure}

\section{Trainable Language Model Experiments}
\subsection{Setup}
We evaluate the method in a 2-layer transformer language model with hidden size 128, four attention heads, sequence length 32, vocabulary size 5000, and routing budget $K=64$. We compare three conditions: a learned top-1 routing baseline, fixed Hopf routing, and randomly permuted fixed routing.

\subsection{Penn Treebank and WikiText-2}
On Penn Treebank, fixed geometric routing reaches within 8\% of the learned-gating baseline in validation perplexity while retaining a concentrated routing footprint and using no learned gate matrix. The same comparison on WikiText-2 yields a nearly identical HOPF/BASELINE ratio under the same training setup.

\begin{table*}[t]
\centering
\caption{Cross-dataset summary.}
\label{tab:crossdataset}
\small
\begin{tabular}{lcc}
\toprule
Metric & PTB & WikiText-2 \\
\midrule
HOPF/BASELINE PPL ratio & 1.081 & 1.081 \\
HOPF vs. PERMUTED $|\Delta\mathrm{ppl}|$ & 0.13 & 0.03 \\
HOPF vs. DENSE eff. ratio & $1.39\times$ & $1.56\times$ \\
\bottomrule
\end{tabular}
\end{table*}

\subsection{Null result}
A key result is that HOPF and PERMUTED are effectively indistinguishable in the trainable language model. This suggests that in the small trainable regime studied here, expert weights co-adapt to the fixed routing scaffold, and the specific geometry does not provide additional quality advantage.

This sharpens the claim of the paper: the sparse fixed-routing mechanism transfers, but the specific Hopf geometry does not show a trainable-LM quality advantage over random fixed routing.

\section{Discussion and Limitations}
The empirical picture is therefore mixed but informative. On the one hand, the proxy experiments show a strong sublinear routing footprint derived from embedding geometry alone. On the other hand, end-to-end language-model training absorbs the geometry-specific advantage, leaving a narrower result: fixed routing can replace learned gating at moderate cost in a toy-scale setting, but the particular geometric scaffold is not uniquely beneficial in that trainable regime.

This work has several important limitations. First, the strongest scaling-law results are established on a semantic proxy rather than large end-to-end language models. Second, the trainable experiments use a 2-layer model and should not be interpreted as evidence for large-scale replacement of learned routing. Third, the Hopf geometry does not outperform randomly permuted fixed routing in the trainable language model setting. Fourth, the experiments do not address attention-side routing. Finally, all experiments were conducted by a single researcher.

\section{Conclusion}
We presented a fixed geometric routing method based on a Hopf-fibration-derived sector map and evaluated it on both a semantic proxy and a small trainable language model. The main empirical result is that angular structure in normalized embeddings supports a sublinear effective routing footprint and can partially transfer to a trainable routing setting without a learned gate matrix. The main limitation is that the specific geometry does not provide an observable quality advantage over random fixed routing in the trainable LM regime studied here.

The contribution is therefore narrow but concrete: embedding angular structure appears relevant not only for compression but also for routing computation.

\bibliographystyle{plainnat}
\bibliography{references}

\end{document}